\section{Conclusiones}
\label{sec:conclusiones}

Este trabajo diseñó e implementó un puente \textit{lock-and-mint} con Ethereum Sepolia como
cadena de origen y Optimism Sepolia como cadena de destino de referencia, en el que el destino
acuña un activo solo tras verificar pruebas de conocimiento cero sobre el bloqueo realizado en el
origen, y comparó dos herramientas para generar esas pruebas, Noir y SP1. La transferencia
completa se ejecutó desde Ethereum Sepolia hacia una red local con las reglas de Optimism
Sepolia, y en esta última red solo se comprobó, mediante llamadas que no modifican el estado, el
verificador de SP1 desplegado. Las conclusiones se organizan según los objetivos específicos
planteados en el Capítulo~\ref{chp:introduccion}.

Respecto del primer objetivo, el estudio de los mecanismos de prueba mostró que la
diferencia entre ambas herramientas no reside en el lenguaje, sino en el objeto que se demuestra.
Noir compila cada programa en un circuito específico, que Barretenberg demuestra con el sistema
UltraHonk sobre una cadena de referencia universal, de modo que cada circuito tiene su propia clave
de verificación y su propio contrato verificador
(Secciones~\ref{sec:noir_flujo} y~\ref{sec:noir_impl_publicas}). El costo de una prueba queda
determinado por el número de restricciones del circuito, que incluye ambas ramas de cada
condicional y las restricciones correspondientes a los tamaños máximos de los datos. SP1, en
cambio, demuestra la ejecución de un programa sobre una máquina virtual RISC-V, cuyo costo depende
de los ciclos efectivamente ejecutados, y comprime la prueba resultante mediante recursión hasta
envolverla en una prueba Groth16 que verifica un contrato común a todos los programas
(Sección~\ref{sec:sp1_flujo}). Esta diferencia de arquitectura anticipa el resultado principal de
las mediciones, según el cual Noir concentra su costo en la preparación del circuito y SP1 en cada
prueba.

El segundo objetivo se cumplió, ya que ambas herramientas permitieron expresar los dos enunciados
que sustentan el puente (Secciones~\ref{sec:noir_impl} y~\ref{sec:sp1_impl}). La prueba de
consenso deriva la raíz de firma del encabezado atestiguado, verifica la firma agregada de los
miembros participantes del comité de sincronización, exige la participación de más de dos tercios
del comité y vincula el bloque de ejecución con el encabezado mediante una rama SSZ. La prueba de
inclusión establece que la transacción de bloqueo y su recibo ocupan la misma posición de los
árboles de transacciones y de recibos de ese bloque y que el recibo registra una ejecución
exitosa. Las implementaciones no son idénticas. En Noir, la prueba de consenso se divide en dos
circuitos porque el circuito único no pudo compilarse en el equipo disponible, el enlace entre
ambas pruebas es el par de raíces del bloque y el circuito de inclusión comprueba, además, el
evento \texttt{Locked} (Sección~\ref{sec:puente_variantes}), sin que esas diferencias alteren las
garantías del puente. Los contratos de destino solo acuñan cuando ambas
pruebas se aceptan y sus valores públicos satisfacen las condiciones de la
Sección~\ref{sec:puente_aceptacion}, de manera que el \textit{relayer} deja de ser un acuñador de
confianza y pasa a ser un probador sin permisos especiales. La implementación en Noir completó,
además, transferencias reales de Ethereum Sepolia a Optimism Sepolia en unos tres minutos y medio
cada una, con un gas de anclaje y de acuñación que coincidió con el medido en la red local. La
firma del comité acredita la autenticidad del encabezado atestiguado, pero no la finalidad del
bloque, y la implementación conserva, por ello, dos supuestos de confianza. El puente no verifica
la finalidad del bloque anclado, y la rotación del comité depende del propietario de
\texttt{BeaconLightClient}, que actualiza el compromiso sin una prueba de la transición.

En cuanto al tercer objetivo, las mediciones caracterizaron el costo de ambas herramientas
en la generación y en la verificación de las pruebas. En la generación de la prueba de consenso,
Noir requirió 89,3 minutos de costo fijo, dominado por una compilación que alcanzó 86,7 GiB de
memoria, y 3,0 minutos por prueba con un máximo de 19,0 GiB, mientras que SP1 requirió un costo
fijo de segundos, medido con el ejecutable ya compilado, y 17,2 minutos por prueba con unos 86 GiB de memoria en cada una
(Secciones~\ref{sec:resultados_etapas} y~\ref{sec:resultados_amortizacion}). Con las etapas
medidas, SP1 obtuvo la primera prueba unas cinco veces más rápido, aunque esa cifra no incluye la
compilación desde cero de su programa invitado y no representa, por lo tanto, el costo completo de
preparación de ambas herramientas. Noir resulta, en cambio, 5,7
veces más rápido por prueba en operación, y su tiempo total acumulado es menor a partir de la
séptima prueba generada con un mismo programa. Estas cifras de Noir corresponden, sin embargo, a una
versión anterior de su circuito de consenso, que demuestra un enunciado menor que el de SP1, aunque
los dos circuitos que demuestran el enunciado completo se probaron en unos 98 segundos en un equipo
de 18 GB de memoria. En la verificación, en cambio, con las reglas vigentes en la cadena de
destino, la prueba Groth16 de SP1 requiere unas 259\,000 unidades de gas, frente a entre 4,49 y
4,75 millones de cada prueba UltraHonk de Noir, cuyo anclaje, con dos pruebas, consume unos
9,08 millones y cuya acuñación, unos 4,87 millones. Los datos de SP1 ocupan, además, entre 14 y
20 veces menos (Sección~\ref{sec:resultados_gas}). El
costo de verificación de Noir depende sobre todo del número de entradas públicas, mientras que el
de SP1 es prácticamente independiente del programa. En cuanto a la complejidad del desarrollo, SP1
permitió escribir los programas en Rust con datos de tamaño variable, ejecutarlos sin generar la
prueba y obtener su número de ciclos en segundos, mientras que en Noir cada modificación de un
circuito de consenso exige repetir una compilación que, en las mediciones, demandó entre 32 minutos
y casi una hora y media, y el programa de inclusión obliga a fijar de antemano los tamaños máximos
de sus datos.

En conjunto, los resultados muestran que, en las condiciones evaluadas, ninguna de las dos
herramientas domina a la otra en todas las dimensiones. Para el programa de consenso medido, Noir
ofrece el menor costo de generación por prueba y el menor requisito de memoria del probador una
vez compilado el circuito, mientras que SP1 ofrece el menor costo de verificación en la cadena de
destino, el menor tamaño de las pruebas y la mayor flexibilidad de desarrollo. La elección entre
ambas depende, por lo tanto, de si el costo
determinante para la operación del puente es el del equipo que genera las pruebas o el del gas que
se paga en cada transacción. Estas conclusiones están sujetas a las condiciones de validez de la
Sección~\ref{sec:metodologia_validez} y a las mediciones que no pudieron realizarse, entre ellas el
tiempo de la envoltura Groth16 de SP1, las etapas de generación del programa de inclusión, la
generación en el servidor de los circuitos de consenso que integran el puente en Noir, el gas de
una transacción completa de anclaje o de acuñación en SP1 y el costo monetario en una red
principal, ya que los importes de Optimism Sepolia corresponden a una red de prueba. Las razones entre ambas herramientas que se mencionan en esta sección describen, por lo
tanto, la configuración evaluada y no una propiedad general de Noir ni de SP1.

\section{Verificación de la hipótesis}
\label{sec:verificacion_hipotesis}

La hipótesis sostenía que el paradigma de máquina virtual de conocimiento cero de propósito
general, representado por SP1, presenta ventajas relevantes en flexibilidad de desarrollo y en
tiempos de ejecución frente a un lenguaje de dominio específico como Noir, al implementar la
verificación de consenso e inclusión de transacciones en Ethereum. Los resultados permiten
evaluarla de forma diferenciada para cada una de sus dos dimensiones, en términos de diferencias
observables en las condiciones evaluadas y no de significancia estadística.

En la flexibilidad de desarrollo, la hipótesis se confirma. SP1 admite datos de tamaño variable,
como el número de miembros del comité o la longitud de las transacciones, sin fijar máximos que
condicionen el circuito, recorre solo las claves de los firmantes efectivos y permite ejecutar el
programa sin generar la prueba para depurarlo y estimar su costo. La derivación de su clave de
verificación demandó segundos, frente a la compilación de entre 32 minutos y casi una hora y media
y a la obtención de la clave que exige cada modificación de un circuito de consenso en Noir. Esta comparación es
cualitativa en lo que respecta a la compilación de SP1, cuya duración no se midió desde cero.

En los tiempos de ejecución, la hipótesis se confirma solo para la primera prueba y se refuta para
la operación del puente. En las etapas medidas, con su ejecutable ya compilado, SP1 obtuvo la
primera prueba unas cinco veces más rápido que Noir, pero
cada prueba posterior demandó 5,7 veces más tiempo y unas cuatro veces más memoria, de modo que
Noir resulta más rápido a partir de la séptima prueba. Como el programa de consenso se mantiene
fijo durante la rotación del comité, la operación del puente supera ese punto con los primeros
anclajes. Esta conclusión corresponde al caso medido, en el que SP1 se ejecutó solo en el
procesador y sin precompilados, condiciones que lo desfavorecen, y en el que tampoco se midió la
envoltura Groth16, que aumentaría su costo por prueba. El circuito de Noir medido demuestra, en
cambio, un enunciado menor que el de SP1, condición que favorece a Noir, aunque la generación de
los circuitos completos en otro equipo sugiere que esa diferencia no altera el sentido del
resultado.

La hipótesis se confirma, por lo tanto, de forma parcial. SP1 presenta ventajas relevantes
en la flexibilidad de desarrollo y en el tiempo medido para obtener la primera prueba, pero no
en el tiempo de generación de las pruebas durante la operación del puente, en el que la ventaja
corresponde a Noir. Entre las dimensiones complementarias, sobre las que la hipótesis no formula
una predicción, el costo de verificación favorece a SP1, cuya verificación requiere entre 17 y 18
veces menos gas que cada verificación de Noir con las reglas vigentes en la cadena de destino,
mientras que el consumo de memoria del probador en operación favorece a Noir.

\section{Trabajos futuros}
\label{sec:trabajos_futuros}

Los resultados y las limitaciones señaladas permiten identificar varias líneas de continuación.
La primera consiste en completar las mediciones. Para SP1, corresponde medir la envoltura Groth16,
la compilación desde cero del programa invitado y el efecto de los precompilados de SHA-256 y de
BLS12-381, de la aceleración por GPU y de la red de probadores~\cite{succinctnetwork} sobre el
costo por prueba y sobre el punto de equilibrio. Para ambas herramientas, corresponde medir las
etapas de generación del programa de inclusión y repetir cada medición para estimar su
variabilidad. En Noir, corresponde medir en el servidor los circuitos de consenso que integran el
puente, y en SP1, el gas de las transacciones completas de anclaje y de acuñación con pruebas reales
de sus programas, incluida la verificación de la prueba de inclusión. El envío de las
transacciones de SP1 a Optimism Sepolia permitiría, además, comparar el gas de la capa 2 y los
datos publicados en la capa 1 de ambas herramientas, y la aplicación de los precios de una red
principal a esas magnitudes, estimar su costo monetario.

La segunda línea consiste en reducir los supuestos de confianza que conserva el puente. La
primera es la verificación de la finalidad. En ambas implementaciones, el programa de consenso
demuestra que el comité firmó el encabezado que contiene el bloque de ejecución, pero no que ese
bloque esté finalizado, de modo que un bloque que la cadena de origen descarte podría originar una
acuñación sin respaldo. El programa debería recibir el encabezado finalizado y la rama SSZ que lo
vincula con el atestiguado (Sección~\ref{sec:eth_cadena_verificacion}), verificarla y anclar el
bloque de ejecución del encabezado finalizado, lo que exigiría exponer su \textit{slot} entre los
valores públicos y añadiría una espera de cerca de 12,8 minutos antes de cada acuñación
(Sección~\ref{sec:eth_consenso}), con un efecto sobre el costo de las pruebas que debería medirse.
La verificación de la transición del comité mediante una prueba, en lugar
de su actualización por el propietario, eliminaría la dependencia de ese actor, como ocurre en los
clientes ligeros basados en pruebas descritos en la Sección~\ref{sec:puentes_estado_arte}. Por
último, la exigencia de una llamada directa a \texttt{Sender} excluye los bloqueos que llegan a ese
contrato mediante una llamada interna, como el realizado desde una cuenta inteligente durante las
pruebas en Optimism Sepolia, que se ejecutó en el origen pero no permitió acuñar. Demostrar la ejecución de la
transacción, y no solo su llamada de nivel superior, o aceptar los bloqueos que llegan mediante la
delegación de la EIP-7702~\cite{eip7702} ampliaría el conjunto de usuarios que pueden emplear el
puente, y un mecanismo de devolución evitaría que los fondos de una operación no demostrable
queden en custodia de forma indefinida.

La tercera línea consiste en reducir el costo de verificación. En el circuito de inclusión de Noir,
la representación de los valores de 32 bytes en dos mitades de 128 bits, en lugar de un elemento
por byte, como ya ocurre con las raíces en el circuito del comité, reduciría el número de
entradas públicas, un factor que incide de forma considerable en su gas
(Sección~\ref{sec:resultados_gas}). En ambas herramientas, la agregación de las pruebas de consenso
y de inclusión en una sola mediante recursión permitiría verificar una única prueba por
acuñación. La disponibilidad en Optimism de los precompilados de BLS12-381 de la propuesta
EIP-2537~\cite{eip2537} abre, por su parte, la posibilidad de comparar la verificación mediante
pruebas con la verificación directa de la firma agregada en el contrato.

Por último, el puente puede extenderse con el flujo de retorno, de quema y liberación, que quedó
fuera del alcance del trabajo, con su despliegue en redes principales y con la incorporación de
otras cadenas de origen o de destino.
